
\textbf{Proposition 2.2.} \uline{Avec les notations et les hypothèses de} 2.1, \uline{soit} @@M0@@ \uline{le morphisme structural de} @@M1@@. \uline{Alors l'homomorphisme} :
@@M2@@
\uline{est surjectif}, \uline{et son noyau est le sous-groupe de} @@M3@@ \uline{engendré par les sous-groupes} @@M4@@ \uline{pour} @@M5@@, \uline{ou}, \uline{ce qui revient au même}, \uline{pour} @@M6@@ \uline{parcourant un système de générateurs de} @@M7@@.

On suppose d'abord @@M8@@, i.e.\ @@M9@@. Alors @@M10@@, et on a pour tout @@M11@@ :
@@M12@@
Par suite, si on pose @@M13@@, on a @@M14@@, ce qui montre que @@M15@@ s'identifie au complémentaire de la section nulle dans le fibré @@M16@@. Considérant @@M17@@ comme un sous-schéma fermé de @@M18@@ par la section nulle, on a d'après 1.1 la suite exacte :
@@M19@@
où @@M20@@ est la section nulle de @@M21@@. Par ailleurs, @@M22@@ se factorise en @@M23@@ où @@M24@@ est le morphisme structural de @@M25@@. D'après 1.6, @@M26@@ est un isomorphisme, inverse de @@M27@@. On a donc la suite exacte :
@@M28@@
d'où l'on tire d'après 1.8 :
@@M29@@
où @@M30@@ est la classe dans @@M31@@ de @@M32@@. Comme pour tout @@M33@@, @@M34@@ est multiple de @@M35@@, ceci prouve 2.2 dans le cas où @@M36@@ est de rang @@M37@@.

